\usepackage{xeCJK}
\usepackage{geometry}
\geometry{left=2cm,right=2cm,top=2.5cm,bottom=2.5cm}
\setlength{\headheight}{12.7pt}

\usepackage{titlesec}
\titleformat{\section}
  {\normalfont\large\bfseries}
  {\thesection}
  {1em}
  {}
\titlespacing*{\section}
  {0pt}
  {3.5ex plus 1ex minus .2ex}
  {2.3ex plus .2ex}

\usepackage{amsmath}
\usepackage{amssymb}
\usepackage{amsthm}
\usepackage{enumitem}
\setlist[enumerate]{itemsep=0pt,topsep=0pt,parsep=0pt,leftmargin=0pt,itemindent=2.5em}
\setlist[itemize]{itemsep=0pt,topsep=0pt,parsep=0pt,leftmargin=0pt,itemindent=2.5em}
\usepackage{fancyhdr}
\pagestyle{fancy}
\fancyhf{}
\usepackage{graphicx}
\usepackage{hyperref}
\usepackage{indentfirst}
\usepackage{lmodern}


\begin{document}
\setlength{\abovedisplayskip}{1pt}
\setlength{\belowdisplayskip}{1pt}

\section{重合引理}

\hypertarget{sec:定理1.1}{}
\noindent\textbf{定理1.1 (项的重合引理)}

给定一阶语言$\mathcal{L}$和
$\mathcal{L}$-\href{../01 一阶语言的结构/01 一阶语言的结构#nameddest=sec:定义1.1}
{结构} $A$, 对任意的项$t$和$A$上的变元赋值$\sigma,\tau$, 如果
\[
    \forall x\in\mathrm{var}(t):x^\sigma=x^\tau,
\]
记作$\sigma\overset{\displaystyle t}{=}\tau$, 那么$t^{(A,\sigma)}=t^{(A,\tau)}$.

\noindent{\kaishu 证明.}

\begin{enumerate}
    \item 对任意的常元$c$, 始终有$c^{(A,\sigma)}=c^A=c^{(A,\tau)}$.
    \item 对任意的变元$x$, 假设$\sigma\overset{\displaystyle x}{=}\tau$,
    因为$x\in\mathrm{var}(x)$, 所以$x^{(A,\sigma)}=x^\sigma=x^\tau=x^{(A,\tau)}$.
    \item 对任意的$n$元函数符号$f$, 假设$t_1,\cdots,t_n$满足命题
    并且$t=ft_1\cdots t_n$满足$\sigma\overset{\displaystyle t}{=}\tau$,
    则对每个$t_i$有
    \vspace{3pt}
    \[
        \mathrm{var}(t_i)\subset\mathrm{var}(t)
        \quad\implies\quad\sigma\overset{\displaystyle t_i}{=}\tau
        \quad\implies\quad t_i^{(A,\sigma)}=t_i^{(A,\tau)},
    \]
    从而$t^{(A,\sigma)}=f^A(t_1^{(A,\sigma)},\cdots,t_n^{(A,\sigma)})
    =f^A(t_1^{(A,\tau)},\cdots,t_n^{(A,\tau)})=t^{(A,\tau)}$. \hfill $\qedsymbol$
\end{enumerate}

\vspace{10pt}

\hypertarget{sec:定理1.2}{}
\noindent\textbf{定理1.2 (公式的重合引理)}

给定一阶语言$\mathcal{L}$和$\mathcal{L}$-结构$A$,
对任意的公式$\varphi$和$A$上的变元赋值$\sigma,\tau$, 如果
\[
    \forall x\in\mathrm{FV}(\varphi):x^\sigma=x^\tau,
\]
记作$\sigma\overset{\displaystyle\varphi}{=}\tau$,
那么$(A,\sigma)\models\varphi$当且仅当$(A,\tau)\models\varphi$.

\noindent{\kaishu 证明.}

\begin{enumerate}
    \item 对任意的项$t_1,t_2$, 假设$\varphi=t_1\equiv t_2$满足
    $\sigma\overset{\displaystyle\varphi}{=}\tau$, 则对每个$t_i$有
    \[
        \mathrm{var}(t_i)\subset\mathrm{FV}(\varphi)
        \quad\implies\quad\sigma\overset{\displaystyle t_i}{=}\tau
        \quad\implies\quad t_i^{(A,\sigma)}=t_i^{(A,\tau)},
    \]
    从而
    \[
        (A,\sigma)\models\varphi\quad\iff\quad
        t_1^{(A,\sigma)}=t_2^{(A,\sigma)}\quad\iff\quad
        t_1^{(A,\tau)}=t_2^{(A,\tau)}\quad\iff\quad(A,\tau)\models\varphi.
    \]

    \item 对任意的$n$元关系符号$r$和项$t_1,\cdots,t_n$,
    假设$\varphi=rt_1\cdots t_n$满足$\sigma\overset{\displaystyle\varphi}{=}\tau$.
    同理每个$t_i$赋值相同, 从而
    \[
        (A,\sigma)\models\varphi\,\iff\,
        (t_1^{(A,\sigma)},\cdots,t_n^{(A,\sigma)})\in r^A\,\iff\,
        (t_1^{(A,\tau)},\cdots,t_n^{(A,\tau)})\in r^A\,\iff\,
        (A,\tau)\models\varphi.
    \]

    \item 假设$\varphi_1$满足命题并且$\varphi=\neg\varphi_1$满足
    $\sigma\overset{\displaystyle\varphi}{=}\tau$.
    因为$\mathrm{FV}(\varphi)=\mathrm{FV}(\varphi_1)$, 所以
    $\sigma\overset{\displaystyle\varphi_1}{=}\tau$, 从而
    \[
        (A,\sigma)\models\varphi\,\iff\,\neg\,(A,\sigma)\models\varphi_1\,\iff\,
        \neg\,(A,\tau)\models\varphi_1\,\iff\,(A,\tau)\models\varphi.
    \]

    \item 假设$\varphi_1,\varphi_2$满足命题并且$\varphi=(\varphi_1\to\varphi_2)$
    满足$\sigma\overset{\displaystyle\varphi}{=}\tau$.
    
    \vspace{3pt}\hspace{2.17em}
    因为$\mathrm{FV}(\varphi)=\mathrm{FV}(\varphi_1)\cup\mathrm{FV}(\varphi_2)$,
    所以$\sigma\overset{\displaystyle\varphi_1}{=}\tau$并且
    $\sigma\overset{\displaystyle\varphi_2}{=}\tau$, 同理
    $(A,\sigma)\models\varphi\iff(A,\tau)\models\varphi$.

    \vspace{3pt}
    \item 对任意的变元$x$, 假设$\varphi_1$满足命题并且$\varphi=\exists x\varphi_1$
    满足$\sigma\overset{\displaystyle\varphi}{=}\tau$.

    \vspace{3pt}\hspace{2.17em}
    假设$(A,\sigma)\models\varphi$, 即存在$a\in A$使得
    $\displaystyle\left(A,\sigma\frac{a}{x}\right)\models\varphi_1$,
    下证$\displaystyle\sigma\frac{a}{x}\overset{\displaystyle\varphi_1}{=}
    \tau\frac{a}{x}$. 对任意的$y\in\mathrm{FV}(\varphi_1)$:

    \vspace{3pt}\hspace{2.17em}
    如果$y=x$, 那么$\displaystyle\sigma\frac{a}{x}(y)=a=\tau\frac{a}{x}(y)$;
    如果$y\neq x$, 那么$y\in\mathrm{FV}(\varphi_1)\setminus\{x\}
    =\mathrm{FV}(\varphi)$, 依假设有
    \vspace{3pt}
    \[
        \sigma\frac{a}{x}(y)=\sigma(y)=\tau(y)=\tau\frac{a}{x}(y).\\[3pt]
    \]
    于是$\displaystyle\sigma\frac{a}{x}\overset{\displaystyle\varphi_1}{=}
    \tau\frac{a}{x}$, 依假设有
    $\displaystyle\left(A,\tau\frac{a}{x}\right)\models\varphi_1$,
    从而$(A,\tau)\models\varphi$. 反之同理. \hfill $\qedsymbol$
\end{enumerate}





\newpage

\section{等价赋值的观点}

\hypertarget{sec:定义2.1}{}
\noindent\textbf{定义2.1}

给定一阶语言$\mathcal{L}$, $\mathcal{L}$-结构$A$和公式$\varphi$. 显然
$\overset{\displaystyle\varphi}{=}$是$^{\mathcal{V}}\!A$上的等价关系,
把其商集记为$^{\mathcal{V}}\!A/\varphi$.

\vspace{10pt}

\hypertarget{sec:定理2.1}{}
\noindent\textbf{定理2.1}

给定一阶语言$\mathcal{L}$, $\mathcal{L}$-结构$A$和公式$\varphi$, 则存在唯一的
双射$r:\,^{\mathcal{V}}\!A/\varphi\to\,^{\mathrm{FV}(\varphi)}A$满足:
\[
    \forall s\in\,\!^{\mathcal{V}}\!A/\varphi,\,\forall\sigma\in s:\,
    r(s)=\sigma\!\restriction_{\mathrm{FV}(\varphi)}.
\]

\noindent{\kaishu 证明.}

依定义, 当$\sigma\overset{\displaystyle\varphi}{=}\tau$时一定有
$\sigma\!\restriction_{\mathrm{FV}(\varphi)}=\tau\!\restriction_{\mathrm{FV}(\varphi)}$,
所以有唯一的商映射$r:\,^{\mathcal{V}}\!A/\varphi\to\,^{\mathrm{FV}(\varphi)}A$
满足命题.

然后证明$r$是单射. 如果$r(s)=r(s')$, 那么取$\sigma\in s$和$\sigma'\in s'$就有
\[
    \sigma\!\restriction_{\mathrm{FV}(\varphi)}=r(s)=r(s')=\sigma'\!\restriction_{\mathrm{FV}(\varphi)},
\]
从而$\sigma\overset{\displaystyle\varphi}{=}\sigma'$, 即$s=s'$.

最后证明$r$是满射. 对任意的$f\in\,\!^{\mathrm{FV}(\varphi)}A$,
取它的一个扩张$f\subset\sigma\in\,\!^{\mathcal{V}}\!A$,
则$r[\sigma]=\sigma\!\restriction_{\mathrm{FV}(\varphi)}=f$. \hfill $\qedsymbol$

\vspace{10pt}

对任意的$s\in\,\!^{\mathcal{V}}\!A/\varphi$, 如果存在$\sigma\in s$使得
$(A,\sigma)\models\varphi$, 就记$(A,s)\models\varphi$.

根据重合引理, 如果$(A,s)\models\varphi$, 那么对任意的$\tau\in s$都有
$(A,\tau)\models\varphi$.

\vspace{10pt}

此时再给定变元$x$和$a\in A$. 当$\sigma\overset{\displaystyle\varphi}{=}\tau$
时一定有$\displaystyle\sigma\frac{a}{x}\overset{\displaystyle\varphi}{=}
\tau\frac{a}{x}$, 即$\displaystyle\left[\sigma\frac{a}{x}\right]=
\left[\tau\frac{a}{x}\right]$. 所以能定义$\displaystyle s\frac{a}{x}$使得:
\vspace{6pt}
\[
    \forall s\in\,\!^{\mathcal{V}}\!A/\varphi,\,\forall\sigma\in s:\,
    s\frac{a}{x}=\left[\sigma\frac{a}{x}\right].
\]

\end{document}